\section{Introduction}
\label{sec:intro}

In this measurement paper, we investigate the interplay of two important aspects of today's cloud systems: (1) the widespread adoption of network prioritization and (2) the prevalent use of TCP as the primary transport protocol.
Our specific focus is on understanding the suitability of TCP for low-priority flows, and we frame our inquiry as: \textit{Is TCP good enough for low-priority flows under a cloud setting?}

Modern cloud systems host a multitude of applications with varying performance objectives.
To support these needs, there has been a widespread move towards adopting network prioritization (i.e., using priority queues inside switches). 
Network prioritization guarantees isolation across different types of applications, enabling various use-cases such as network flow scheduling (e.g., DAS~\cite{das}, pFabric~\cite{pfabric}, Homa~\cite{homa}, Baraat~\cite{baraat}) and workload co-location (e.g., perfIso~\cite{perfIso}, Heracles~\cite{heracles}, SmartHarvest~\cite{smart-harvest}).
For example, EBB~\cite{ebb}, Meta's backbone network, services several traffic classes in a strictly prioritized manner.
A new use-case of prioritization is in the context of RDMA-based stacks where network operators can use the TCP/IP traffic as a reliable and robust low-priority backup, with RDMA-based traffic as the higher priority traffic in the system~\cite{pangu}.
For many of the above use-cases, the performance of low-priority flows -- albeit less important than the high priority traffic -- is still important. 
For example, in systems like DAS~\cite{das} which uses duplication to reduce tail latency, a lower priority flow may determine the application performance if the corresponding high priority flow becomes a straggler. 

On the other hand, despite a large body of work highlighting the performance implications of using TCP for latency-critical applications~\cite{pdq,p3,pfabric,d2tcp,dctcp}, it remains the defacto protocol in today's cloud systems for important reasons such as backward compatibility and stability~\cite{morgan-stanley-nsdi15,accelnet,pangu,luna}. 
For instance, a recent work from Alibaba~\cite{luna} emphasized that their preference for TCP stems from the lack of specialized in-network support and the presence of a fleet of servers lacking support for emerging technologies (e.g., DPDK).

This motivates an investigation into potential issues that may arise when TCP is used for low-priority flows under network prioritization.
To the best of our knowledge, there has been no prior research dedicated to investigating this specific inquiry.
Intuitively, the use of TCP for low-priority flows can introduce a range of issues. This is primarily because low-priority flows are susceptible to extended queuing delays, which, in the worst-case scenario, can be unbounded~\cite{conquest}. These delays impact TCP's feedback loop, potentially resulting in unnecessary backoffs, and spurious retransmissions, where packets are mistakenly assumed to be lost while they are stuck within network switches.

Building on this intuition, our study evaluates the performance of TCP for low-priority flows across various important use-cases, encompassing a wide range of settings, including characteristics of high and low-priority traffic, the load and workload profiles, and specific TCP and switch configurations (e.g., RTO\textsubscript{min}, buffer size, etc.). 
Our study has two primary objectives: First, to shed light on the cases where TCP is a suitable choice for low-priority flows and identify scenarios where it is not.
Second, to explore minor adjustments that can help address any limitations observed in TCP's use for low-priority flows. Our methodology involves conducting experiments on a small-scale testbed and NS3 simulations~\cite{ns3}. We use a subset of experiments to cross-validate the results on both these platforms and then use simulations to explore a wide range of scenarios.

To quantify the impact of network prioritization and to uncover the room for improving performance of low-priority flows beyond what TCP offers, we consider a \emph{near-optimal} (\nearopt{}) fair transport protocol as a baseline. The \nearopt{} protocol has the advantage of instantaneous access to network state information, including available bandwidth, queue lengths, and packet loss details. 

Consequently, it represents an upper-bound that an efficient \emph{fair} transport protocol can achieve for low-priority flows. 
Our key measurement insights are as follows:
\begin{itemize}
    \item For several use-cases (e.g., network scheduling), TCP's performance for low-priority flows is within $2\times{}$ of \nearopt{}. This is because, for such use-cases, low-priority flows frequently get some fraction of network share allowing TCP to continually adapt to the network conditions.

    \item Low-priority TCP flows suffer the most --- more than $12\times{}$ compared to \nearopt{} --- when they coexist with high priority traffic that exhibits an \emph{on-off} behavior (typical of distributed ML model training workloads~\cite{p3,tiresias}). The \emph{on} period can completely starve the low-priority queue, causing TCP's feedback loop to stall and experience severe performance degradation. During the \emph{off} period, unnecessary congestion backoffs limit TCP from efficiently utilizing the link bandwidth. 

    \item Emerging technology trends such as increased link speeds and relatively smaller switch buffer sizes~\cite{one-more-config} significantly affect TCP's performance. Specifically when available buffer is equally divided between all the queues, TCP's performance is degraded by up to $5\times{}$ compared to when the entire buffer is available for the low-priority queue. 

    \item Two simple strategies -- weighted fair queuing (WFQ), and \emph{cross-queue} congestion notification -- effectively alleviate stalls in TCP's feedback loop, resulting in up to 9$\times$ reduction in completion times for low-priority flows.
        
\end{itemize}

The paper is organized as follows: We first present popular and emerging use-cases of network prioritization and highlight their performance implications for TCP when used for low-priority flows (\S\ref{sec:motive}). In~\S\ref{sec:eval-main}, we measure TCP's performance for low-priority flows across popular use-cases.
In~\S\ref{sec:micro}, we investigate the role of several configurations crucial to the performance of TCP such as emerging technology trends (e.g., smaller buffer sizes) and important network and transport parameters (e.g., RTO\textsubscript{min}).
Finally, in \S\ref{sec:future}, we discuss two simple strategies (and do the preliminary evaluation) to show substantial improvement in TCP's performance for low-priority flows.